\scanpage{159}
Comme $\mathrm{A}$ est noethérien, il en est de même de $\mathrm{A}'=\mathrm{A}\{\mathrm{T}_1,\ldots,\mathrm{T}_r\}$ (7.5.4), donc \textit{c)} implique que $\mathrm{B}$ est noethérien\,; comme $\unsure{\mathfrak{J}|\mathfrak{I}}'=\mathfrak{J}\mathrm{A}'$ est un voisinage ouvert de $0$ dans $\mathrm{A}'$ tel que les ${\mathfrak{J}'}^{n}$ forment un système fondamental de voisinages de $0$, les images $\mathfrak{J}^{n}\mathrm{B}$ des ${\mathfrak{J}'}^{n}$ forment un système fondamental de voisinages de $0$ dans $\mathrm{B}$, et comme on sait que $\mathrm{B}$ est séparé et complet, $\mathrm{B}$ est un anneau $\mathrm{J}\mathrm{B}$-adique. Enfin, $\mathrm{B}/\mathfrak{J}\mathrm{B}$ est une algèbre (sur $\mathrm{A}/\mathfrak{J}$) quotient de $\mathrm{A}'/\mathfrak{J}\mathrm{A}'=(\mathrm{A}/\mathfrak{J})[\mathrm{T}_1,\ldots,\mathrm{T}_r]$, donc est de type fini, ce qui achève de prouver que \textit{c)} entraîne \textit{a)}.

Si $\mathrm{B}$ est $\mathfrak{J}\mathrm{B}$-adique et noethérien, $\mathrm{B}$ est isomorphe à $\varprojlim \mathrm{B}_n$, où $\mathrm{B}_n=\mathrm{B}/\mathfrak{J}^{n+1}\mathrm{B}$ (7.2.11), et $\mathfrak{J}\mathrm{B}/\mathfrak{J}^2\mathrm{B}$ est un module de type fini sur $\mathrm{B}/\mathfrak{J}\mathrm{B}$. Soit $(a_j)_{1\leqslant j\leqslant s}$ un système de générateurs du $(\mathrm{B}/\mathfrak{J}\mathrm{B})$-module $\mathfrak{J}\mathrm{B}/\mathfrak{J}^2\mathrm{B}$, et $(c_i)_{1\leqslant i\leqslant r}$ un système d’éléments de $\mathrm{B}/\mathfrak{J}^2\mathrm{B}$ dont les classes mod. $\mathfrak{J}\mathrm{B}/\mathfrak{J}^2\mathrm{B}$ forment un système de générateurs de la $(\mathrm{A}/\mathfrak{J})$-algèbre $\mathrm{B}/\mathfrak{J}\mathrm{B}$\,; on voit aussitôt que les $c_i a_j$ forment un système de générateurs de la $(\mathrm{A}/\mathfrak{J}^2)$-algèbre $\mathrm{B}/\mathfrak{J}^2\mathrm{B}$, donc \textit{a)} implique \textit{b)}.

Reste à prouver que \textit{b)} entraîne \textit{c)}. L’hypothèse entraîne que $\mathrm{B}_1$ est un anneau noethérien, et comme $\mathrm{B}_1=\mathrm{B}_2/\mathfrak{J}^2\mathrm{B}_2$, on a $\mathfrak{J}^2\mathrm{B}_1=0$, donc $\mathfrak{J}\mathrm{B}_1=\mathfrak{J}\mathrm{B}_1/\mathfrak{J}^2\mathrm{B}_1$ est un $\mathrm{B}_0$-module de type fini. Les conditions de (7.2.7) sont donc vérifiées par le système projectif $(\mathrm{B}_n)$ et $\mathrm{B}$ est un anneau $\mathfrak{J}\mathrm{B}$-adique. Soit $(c_i)_{1\leqslant i\leqslant r}$ un système fini d’éléments de $\mathrm{B}$ dont les classes mod. $\mathfrak{J}^2\mathrm{B}$ engendrent la $(\mathrm{A}/\mathfrak{J}^2)$-algèbre $\mathrm{B}/\mathfrak{J}^2\mathrm{B}$ (isomorphe à $\mathrm{B}_1$).
% slanted
\margin{\handwritten{définition insuffisante il faut aussi que les $\unsure{c_i|c'}\in J B$ engendrent $JB/J^2B$ !}}
Il existe donc un $\mathrm{A}$-homomorphisme continu $u$ de $\mathrm{A}'=\mathrm{A}\{\mathrm{T}_1,\ldots,\mathrm{T}_r\}$ dans $\mathrm{B}$ qui se réduit à $\varphi$ dans $\mathrm{A}$ et est tel que $u(\mathrm{T}_i)=c_i$ pour $1\leqslant i\leqslant r$ (7.5.3)\,; si nous prouvons que $u$ est \textit{surjectif}, \textit{c)} sera établi, car de $u(\mathrm{A}')=\mathrm{B}$ on déduira $u(\mathfrak{J}^n\mathrm{A}')=\mathfrak{J}^n\mathrm{B}$, autrement dit $u$ sera un morphisme strict d’anneaux topologiques et $\mathrm{B}$ sera donc isomorphe à un quotient de $\mathrm{A}'$ par un idéal fermé. Or, comme $\mathrm{B}$ est complet pour la topologie $\mathfrak{J}\mathrm{B}$-adique, il suffit ([1], p. 18-07) de montrer que l’homomorphisme $\operatorname{grad}(\mathrm{A}')\to\operatorname{grad}(\mathrm{B})$ déduit canoniquement de $u$ pour les filtrations $\mathfrak{J}$-adiques sur $\mathrm{A}'$ et $\mathrm{B}$, est surjectif. Mais par définition, les homomor\uline{phismes $\mathrm{A}'/\mathfrak{J}\mathrm{A}'\to \mathrm{B}/\mathfrak{J}\mathrm{B}$ et $\mathfrak{J}\mathrm{A}'/\mathfrak{J}^2\mathrm{A}'\to \mathfrak{J}\mathrm{B}/\mathfrak{J}^2\mathrm{B}$ déduits de $u$ sont surjectifs\,;} \margin{\handwritten{!}} par récurrence sur $n$, on en déduit aussitôt qu’il en est de même de $\mathfrak{J}\mathrm{A}'/\mathfrak{J}^n\mathrm{A}'\to \mathfrak{J}\mathrm{B}/\mathfrak{J}^n\mathrm{B}$, et \textit{a fortiori} de $\mathfrak{J}^n\mathrm{A}'/\mathfrak{J}^{n+1}\mathrm{A}'\to \mathfrak{J}^n\mathrm{B}/\mathfrak{J}^{n+1}\mathrm{B}$, ce qui achève la démonstration.

\textbf{7.6.} \textit{Anneaux complets de fractions.}

\textbf{(7.6.1)} Soient $\mathrm{A}$ un anneau linéairement topologisé, $(\mathfrak{J}_\lambda)$ un système fondamental de voisinages de $0$ dans $\mathrm{A}$ formé d’idéaux, $\mathrm{S}$ une partie multiplicative de $\mathrm{A}$. Soit $u_\lambda$ l’homomorphisme canonique $\mathrm{A}\to \mathrm{A}_\lambda=\mathrm{A}/\mathfrak{J}_\lambda$, et pour $\mathfrak{J}_\mu\subset \mathfrak{J}_\lambda$, soit $u_{\lambda\mu}$ l’homomorphisme canonique $\mathrm{A}_\mu\to \mathrm{A}_\lambda$. Posons $\mathrm{S}_\lambda=u_\lambda(\mathrm{S})$, de sorte que $u_{\lambda\mu}(\mathrm{S}_\mu)=\mathrm{S}_\lambda$. Les $u_{\lambda\mu}$ donnent canoniquement des homomorphismes surjectifs $\mathrm{S}_\mu^{-1}\mathrm{A}_\mu\to \mathrm{S}_\lambda^{-1}\mathrm{A}_\lambda$, pour lesquels ces anneaux forment un système projectif\,; désignons par $\mathrm{A}\{\mathrm{S}^{-1}\}$ la limite projective de ce système. Cette définition ne dépend pas du système fondamental de voisinages $(\mathfrak{J}_\lambda)$ choisi\,; en effet :

\textit{Proposition} \textbf{(7.6.2)}. --- \textit{L’anneau $\mathrm{A}\{\mathrm{S}^{-1}\}$ est topologiquement isomorphe au séparé complété de l’anneau $\mathrm{S}^{-1}\mathrm{A}$ pour la topologie dont un système fondamental de voisinages de $0$ est formé des $\mathrm{S}^{-1}\mathfrak{J}_\lambda$.}
% archive sheet number: 158
